\chapter{Wann teilen und wann nicht}\label{kap:teilen}

\begin{merksatz}
\boxtitel{Die vier Gründe, die tragen}
\begin{center}
\renewcommand{\arraystretch}{1.5}
\begin{tabular}{@{}P{3.2cm}P{2.0cm}P{2.6cm}@{}}
\toprule
\textbf{Grund} & \textbf{Trägt?} & \\
\midrule
\textbf{Andere Betriebsanforderung} & \textcolor{gruen}{\textbf{ja}} & \textbf{Zeitschranke, Speicher}\\
\textbf{Andere Sprache nötig} & \textcolor{gruen}{\textbf{ja}} & \textcolor{gruen}{Band 21, Kap.~14}\\
\textbf{Eigenständig erneuerbar} & \textcolor{gruen}{\textbf{ja}} & \\
\textbf{Eigenes Team, eigener Takt} & \textcolor{gruen}{\textbf{ja}} & \textcolor{gruen}{Conway}\\
\bottomrule
\end{tabular}
\end{center}

\begin{center}
\begin{minipage}{0.88\textwidth}
\kursiv{Die ersten beiden sind technisch zwingend --- der Messdatenerfasser \kursiv{muss} ein eigener Prozess in einer anderen Sprache sein. Die letzten beiden sind organisatorisch und deshalb ehrlicher, als sie klingen: Zuschnitt nach Zuständigkeit ist ein legitimer Grund, wenn man ihn benennt.}
\end{minipage}
\end{center}

\textbf{Die Gründe, die nicht tragen.}

\begin{center}
{\small
\renewcommand{\arraystretch}{1.5}
\begin{tabular}{@{}P{3.2cm}P{2.0cm}P{2.6cm}@{}}
\toprule
\textbf{Behauptung} & \textbf{Trägt?} & \\
\midrule
\enquote{Besser skalierbar} & \textcolor{mittel}{\textbf{selten}} & \textbf{erst messen}\\
\enquote{Sauberer getrennt} & \textcolor{akzent}{\textbf{nein}} & \textcolor{akzent}{\textbf{Module genügen}}\\
\enquote{Leichter zu prüfen} & \textcolor{akzent}{nein} & \textcolor{akzent}{schwerer}\\
\enquote{Modern} & \textcolor{akzent}{nein} & \\
\bottomrule
\end{tabular}}
\end{center}

\begin{center}
\kursiv{Die zweite Zeile ist die häufigste und die schwächste: Wer eine saubere Trennung will, führt ein Modul ein --- Band 17, Kapitel~4. Ein eigener Prozess erzwingt die Trennung, und das ist der einzige Unterschied. Er kostet achtzehn Personentage.}
\end{center}

\begin{technik}
\boxtitel{Was ein zusätzlicher Bestandteil kostet}
\begin{center}
{\small
\renewcommand{\arraystretch}{1.45}
\begin{tabular}{@{}P{3.0cm}P{2.2cm}P{2.6cm}@{}}
\toprule
\textbf{Posten} & \textbf{Aufwand} & \\
\midrule
Naht entwerfen und bauen & $\approx18$ PT & \textbf{Kap.~\ref{kap:werk}}\\
Bereitstellung, Überwachung & $\approx6$ PT & \textcolor{akzent}{laufend}\\
Ausfallverhalten & $\approx5$ PT & \\
Prüfung über die Naht & $\approx4$ PT & \textcolor{akzent}{\textbf{je Änderung}}\\
Verfügbarkeit sinkt & --- & \textcolor{akzent}{\textbf{multiplikativ}}\\
Verzögerung steigt & $+3$--$80$ ms & \\
\bottomrule
\end{tabular}}
\end{center}

\begin{center}
\begin{minipage}{0.86\textwidth}
\kursiv{Dreiunddreißig Personentage und eine schlechtere Verfügbarkeit --- das ist der Einstandspreis. Er lohnt, wenn einer der vier Gründe vorliegt, und nicht, wenn man nur eine Grenze ziehen will.}
\end{minipage}
\end{center}
\end{technik}
\end{merksatz}

\begin{merksatz}
\boxtitel{Der Zuschnitt}
\begin{center}
{\small
\renewcommand{\arraystretch}{1.5}
\begin{tabular}{@{}P{3.0cm}P{2.4cm}P{2.6cm}@{}}
\toprule
\textbf{Geschnitten nach} & \textbf{Folge} & \\
\midrule
Technischer Schicht & \textcolor{akzent}{\textbf{jede Änderung trifft alle}} & \textcolor{akzent}{\textbf{schlecht}}\\
Fachlichem Bereich & \textcolor{gruen}{\textbf{Änderung bleibt lokal}} & \textcolor{gruen}{gut}\\
Betriebs\-anforderung & \textcolor{gruen}{ja} & \textcolor{gruen}{zwingend}\\
Abteilungsgrenze & \textcolor{mittel}{ehrlich benennen} & \textbf{Conway}\\
\bottomrule
\end{tabular}}
\end{center}

\begin{center}
\begin{minipage}{0.88\textwidth}
\kursiv{Die erste Zeile ist der klassische Fehlschnitt: ein Dienst für die Oberfläche, einer für die Logik, einer für die Daten. Eine neue Befundart ändert dann alle drei --- und jede Änderung braucht drei abgestimmte Bereitstellungen.}
\end{minipage}
\end{center}

\textbf{Die Probe.}

\begin{center}
\renewcommand{\arraystretch}{1.45}
\begin{tabular}{@{}P{3.0cm}P{2.2cm}P{2.8cm}@{}}
\toprule
\textbf{Typische Änderung} & \textbf{Trifft} & \\
\midrule
Neue Befundart & $\mathbf{1}$ & \textcolor{gruen}{\textbf{guter Schnitt}}\\
Neue Warnungsregel & $1$ & \textcolor{gruen}{gut}\\
Neues Feld im Befund & $\mathbf{3}$ & \textcolor{mittel}{unvermeidlich}\\
Andere Anzeige & $1$ & \textcolor{gruen}{gut}\\
\bottomrule
\end{tabular}
\end{center}

\begin{center}
\kursiv{Die Probe ist einfach und wird selten gemacht: Man nimmt die zehn letzten Änderungen und zählt, wie viele Bestandteile jede berührt hat. Liegt der Schnitt richtig, ist der Wert meist eins.}
\end{center}

\begin{warnung}
\boxtitelr{Der Weg zurück}
\begin{center}
\renewcommand{\arraystretch}{1.5}
\begin{tabular}{@{}P{2.8cm}P{2.2cm}P{3.0cm}@{}}
\toprule
& \textbf{Teilen} & \textbf{Zusammen\-legen}\\
\midrule
Aufwand & $\approx33$ PT & \textcolor{gruen}{$\approx\mathbf{8}$ \textbf{PT}}\\
Risiko & \textcolor{mittel}{mittel} & \textcolor{gruen}{gering}\\
Wird gemacht & \textcolor{gruen}{oft} & \textcolor{akzent}{\textbf{fast nie}}\\
\bottomrule
\end{tabular}
\end{center}

\begin{center}
\begin{minipage}{0.88\textwidth}
\kursiv{Zusammenlegen ist viermal billiger als Teilen und wird trotzdem kaum je erwogen --- weil es wie ein Rückschritt aussieht. Dabei ist es die richtige Antwort, sobald zwei Bestandteile immer gemeinsam geändert und gemeinsam bereitgestellt werden.}
\end{minipage}
\end{center}

\begin{center}
\kursiv{Die Prüffrage lautet: Wurden diese beiden in den letzten zwölf Monaten je einzeln bereitgestellt? Lautet die Antwort nein, sind sie ein Bestandteil mit einer überflüssigen Naht dazwischen.}
\end{center}
\end{warnung}
\end{merksatz}

\begin{aufgabe}
\begin{enumerate}[label=(\alph*)]
\item Bewerten Sie das Vorhaben.
\item Rechnen Sie den Einstiegsfall.
\item Bestimmen Sie, womit man anfängt.
\end{enumerate}

\loesung
\textbf{(a) Das Vorhaben bewerten.}

\begin{center}
{\small
\renewcommand{\arraystretch}{1.45}
\begin{tabular}{@{}P{2.4cm}P{2.2cm}P{2.2cm}P{1.8cm}@{}}
\toprule
\textbf{Bestandteil} & \textbf{Grund} & \textbf{Art} & \textbf{Urteil}\\
\midrule
Erfassung & Zeitschranke & \textbf{technisch} & \textcolor{gruen}{\textbf{ja}}\\
Übertragung & \textbf{Verfügbarkeit} & technisch & \textcolor{gruen}{ja}\\
Speicherung & Standard\-system & technisch & \textcolor{gruen}{ja}\\
Auswertung & Fachlogik & \textbf{fachlich} & \textcolor{gruen}{ja}\\
Oberfläche & \textbf{eigener Takt} & organisat. & \textcolor{gruen}{ja}\\
\midrule
\textit{Verlaufsdienst} & \textcolor{akzent}{\enquote{sauberer}} & \textcolor{akzent}{keiner} & \textcolor{akzent}{\textbf{nein}}\\
\bottomrule
\end{tabular}}
\end{center}

\begin{center}
\begin{minipage}{0.88\textwidth}
\kursiv{Fünf Bestandteile mit tragenden Gründen und einer ohne. Der letzte ist der Vorschlag aus Kapitel~\ref{kap:werk} --- er hätte zwei zusätzliche Nähte, dreiunddreißig Personentage und eine schlechtere Verfügbarkeit gekostet, für einen Teil, der sich nie unabhängig ändert.}
\end{minipage}
\end{center}

\medskip
\textbf{(b) Der Einstiegsfall.}

\textbf{Die Lage.} Ein kleineres Haus, ein Team von vier Personen.

\begin{center}
{\small
\renewcommand{\arraystretch}{1.45}
\begin{tabular}{@{}P{2.6cm}P{2.2cm}P{2.2cm}P{1.8cm}@{}}
\toprule
& \textbf{Fünf Teile} & \textbf{Zwei Teile} & \\
\midrule
Erstellung & $\mathbf{34}$ PM & $\mathbf{21}$ PM & \\
Nähte & $4$ & $\mathbf{1}$ & \\
Verfügbarkeit & $99{,}3\,\%$ & $\mathbf{99{,}8\,\%}$ & \\
Bereitstellungen & $5$ & $2$ & \\
Betriebsaufwand & $\approx0{,}8$ Stellen & $\approx\mathbf{0{,}3}$ & \\
\bottomrule
\end{tabular}}
\end{center}

\begin{center}
\begin{minipage}{0.88\textwidth}
\kursiv{Bei vier Personen sind zwei Bestandteile richtig: der Messdatenerfasser, weil die Zeitschranke es erzwingt, und alles andere in einem Programm. Dreizehn Personenmonate weniger und bessere Verfügbarkeit.}
\end{minipage}
\end{center}

\begin{center}
\kursiv{Die Trennung nach Fachbereichen bleibt --- als Module innerhalb eines Programms, mit denselben Regeln aus Band 17. Wenn ein Teil später eigene Betriebsanforderungen bekommt, lässt er sich herauslösen. Der umgekehrte Weg ist teurer.}
\end{center}

\medskip
\textbf{(c) Womit man anfängt.}

\begin{center}
{\small
\renewcommand{\arraystretch}{1.5}
\begin{tabular}{@{}P{1.0cm}P{3.4cm}P{3.4cm}@{}}
\toprule
$1$ & \textbf{So wenige Bestandteile wie möglich} & \textbf{ein Programm, gut zerlegt}\\
$2$ & Herauslösen, was muss & Zeitschranke, Sprache\\
$3$ & \textbf{Messen, wo es klemmt} & \textcolor{gruen}{Kap.~\ref{kap:messen}}\\
$4$ & \textbf{Herauslösen, was sich belegen lässt} & \\
$5$ & Zusammenlegen, was nicht trägt & \textcolor{gruen}{\textbf{regelmäßig prüfen}}\\
\bottomrule
\end{tabular}}
\end{center}

\begin{center}
\begin{minipage}{0.88\textwidth}
\kursiv{Der erste Schritt ist der unbeliebteste und der richtige: mit einem gut zerlegten Programm beginnen und nur herauslösen, was sich begründen lässt. Der umgekehrte Weg --- viele Bestandteile von Anfang an --- legt Nähte fest, bevor man weiß, wo sie hingehören.}
\end{minipage}
\end{center}

\begin{center}
\kursiv{Der fünfte Schritt fehlt in fast allen Vorhaben. Er kostet einen halben Tag im Quartal: die Änderungen der letzten Monate ansehen und fragen, welche Bestandteile immer gemeinsam betroffen waren.}
\end{center}

\begin{warnung}
\boxtitelr{Die Zahl der Nähte wächst schneller als die Zahl der Teile}
\begin{center}
\renewcommand{\arraystretch}{1.45}
\begin{tabular}{@{}P{2.4cm}P{2.2cm}P{2.4cm}P{2.0cm}@{}}
\toprule
\textbf{Teile} & \textbf{Nähte} & \textbf{Aufwand} & \textbf{Verfügbar}\\
\midrule
$2$ & $\mathbf{1}$ & $33$ PT & $99{,}8\,\%$\\
$5$ & $4$ & $132$ PT & $99{,}3\,\%$\\
$12$ & $\mathbf{18}$ & \textcolor{akzent}{$594$ PT} & \textcolor{akzent}{$97{,}9\,\%$}\\
$30$ & \textcolor{akzent}{$\mathbf{60}$} & \textcolor{akzent}{$1980$ PT} & \textcolor{akzent}{$\mathbf{94{,}9\,\%}$}\\
\bottomrule
\end{tabular}
\end{center}

\begin{center}
\begin{minipage}{0.88\textwidth}
\kursiv{Bei dreißig Bestandteilen ist die Verfügbarkeit auf knapp fünfundneunzig Prozent gefallen --- achtzehn Tage Ausfall im Jahr, ohne dass ein einziger Bestandteil schlecht wäre. Das ist die Rechnung aus Kapitel~\ref{kap:messen}, hier als Entscheidungsgrundlage.}
\end{minipage}
\end{center}

\begin{center}
\kursiv{Wer viele Bestandteile will, muss jeden einzelnen verfügbarer machen --- und das kostet je Neun etwa das Dreifache. Der Aufwand wächst damit doppelt: mehr Nähte und höhere Anforderungen an jedes Teil.}
\end{center}
\end{warnung}
\end{aufgabe}
\vspace{-4pt}
\noindent Newman\zit{17} behandelt Zu\-schnitt und Grö\-ße ver\-teil\-ter Diens\-te; Conway\zit{6} erklärt, warum der Schnitt den Zu\-stän\-dig\-kei\-ten folgt; Richardson\zit{21} sammelt die Bau\-for\-men mit ihren Kos\-ten.
